\section{Conclusion}\label{sec:conclusion}

We have developed a focused state-forcing application in which local objecthood and gluing failure can be studied without leaving a conservative extension of ordinary pointwise semantics. The paper does not attempt to reconstruct the full larger logic architecture from which the state-forcing fragment originates; instead, it isolates exactly the components needed for a cohomological analysis of visible and invisible obstruction data. Within that scope, the first structural point is categorical: the set-valued local object functor should not be asked to carry cohomological data by itself. Instead, one passes to a realization prestack, whose stackification has connected-component sheaf equal to the sheafification of the original local object functor. Visible value classes are therefore not an auxiliary notion external to the gerbe theory; they are exactly the global connected components of the realization stack.

This changes the shape of the obstruction theory. In general, an admitted reference does not determine a single gerbe, but a family of component gerbes indexed by visible branches. Gluing-level absence occurs exactly when compatible local data exist but every visible branch gerbe is non-neutral. Thus a single obstruction class is the exceptional case in which the visible branch is unique. Once one fixes a branch, the strict visible quotient, the intrinsic visible quotient, and the homological evaluation map recover the visible part of that branch obstruction. The decisive identification is that intrinsic visibility is controlled by the image of the universal coefficient map on $H_2$, while the blind residue lies in the $\operatorname{Ext}^1(H_1,A)$-summand. This is why nontrivial branch obstructions can survive even when all character channels are admissible.

The multi-branch theory adds a second layer of structure. If the branch label is available, total visibility is governed by the intersection $\bigcap_v K_v$; if one insists on a single channel admissible on every branch, it is governed by the sum $\sum_v K_v$; and the exact branch-preserving strictification budget is $\max_v |K_v|$. Along refinement, admissible branchwise character channels can only increase and hidden branch budgets can only decrease. The point is not a new universal coefficient theorem or a new duality theorem, but a precise semantic identification: in this state-forcing setting, visible information is the homologically detected part of each branch obstruction, while the genuinely blind remainder is the residual $\operatorname{Ext}$-part. The natural next steps are proof-theoretic formulations of state forcing and extensions of the visibility theory beyond the finite abelian setting.
